\documentclass[11pt]{article}

\usepackage{amsmath,amssymb,amsthm}
\usepackage{hyperref}
\usepackage{enumitem}
\usepackage{mathtools}

\title{Cycle--Local Rigidity in Bounded-Degree Graphs}
\author{}
\date{}

\newtheorem{theorem}{Theorem}
\newtheorem{lemma}{Lemma}
\newtheorem{definition}{Definition}

\begin{document}

\maketitle

\section{Introduction}

We prove a structural rigidity theorem showing that graphs whose
radius-$R$ neighborhoods are indistinguishable by $\mathrm{FO}^k$
formulas cannot support arbitrarily large cycle-overlap rank.

The result provides the core combinatorial mechanism used in the
Oblivion Atom program and establishes a bounded cycle complexity
for $\mathrm{FO}^k$-homogeneous bounded-degree graphs.

\section{Setting}

Let $G$ be a finite graph with maximum degree at most $\Delta$.

For a vertex $v$ define the radius-$R$ neighborhood

\[
B_R(v)=\{u\in V(G): d_G(u,v)\le R\}.
\]

Two vertices $u,v$ are said to have the same $\mathrm{FO}^k$ type if

\[
\operatorname{tp}_{FO^k}(B_R(u))
=
\operatorname{tp}_{FO^k}(B_R(v)).
\]

\section{Cycle--Local Rigidity}

\begin{theorem}[Cycle--Local Rigidity]
Let $k,\Delta,R\in\mathbb N$.

There exists a constant $C=C(k,\Delta,R)$ such that for every graph
$G$ with maximum degree $\le\Delta$ satisfying

\[
\forall u,v\in V(G):
\operatorname{tp}_{FO^k}(B_R(u))
=
\operatorname{tp}_{FO^k}(B_R(v)),
\]

the cycle--overlap rank satisfies

\[
\mathrm{COR}(G)\le C.
\]

\end{theorem}

\section{FO$^k$ Type Finiteness}

\begin{lemma}
For fixed $k,\Delta,R$ the number of possible
$\mathrm{FO}^k$ types of rooted radius-$R$ neighborhoods is finite.
\end{lemma}

\begin{proof}

There are finitely many rooted graphs of radius $R$ and
degree $\le\Delta$.
FO$^k$ equivalence partitions this finite set.

\end{proof}

Let

\[
T=T(k,\Delta,R)
\]

denote the number of such types.

\section{Cycle Pumping}

\begin{lemma}[Cycle Pumping]

There exists

\[
L=T(2R+1)
\]

such that every simple cycle $C$ with $|C|>L$ contains two vertices
with identical $\mathrm{FO}^k$ neighborhood type whose connecting
segment can be replaced by a shorter path preserving local structure.

\end{lemma}

\begin{proof}

Along the cycle consider the sequence

\[
\tau_i=\operatorname{tp}_{FO^k}(B_R(c_i)).
\]

Since there are only $T$ types,
two vertices share the same type once the cycle length exceeds
$T(2R+1)$.

The Ehrenfeucht--Fra\"iss\'e strategy between their neighborhoods
allows the intermediate segment to be replaced while preserving
local FO$^k$ structure.

\end{proof}

\section{Vertex Signatures}

For a vertex $v$ on a cycle $C$ define

\[
\sigma(v)=
(\operatorname{tp}_{FO^k}(B_R(v)),\ C\cap B_R(v)).
\]

\begin{lemma}

The number of vertex signatures satisfies

\[
S\le T\cdot 2^{\Delta^{R+1}}.
\]

\end{lemma}

\begin{proof}

The ball $B_R(v)$ contains at most $\Delta^{R+1}$ vertices,
hence $2^{\Delta^{R+1}}$ subsets.

\end{proof}

\section{Cycle Signatures}

For a cycle $C=(v_0,\dots,v_{\ell-1})$ define

\[
\Sigma(C)=(\sigma(v_0),\dots,\sigma(v_{\ell-1})).
\]

\begin{lemma}

The number of cycle signatures satisfies

\[
|\Sigma|\le S^L.
\]

\end{lemma}

\section{Local EF Extension}

Let

\[
r=\min\{R/2,3^k\}.
\]

\begin{lemma}

If

\[
\operatorname{tp}_{FO^k}(B_R(u))
=
\operatorname{tp}_{FO^k}(B_R(v))
\]

then for every subset $S\subseteq B_r(u)$ with $|S|\le k$
there exists a bijection

\[
\phi:S\to S'\subseteq B_r(v)
\]

preserving adjacency relations among vertices of $S$.

\end{lemma}

\begin{proof}

This follows from the $k$-round Ehrenfeucht--Fra\"iss\'e game
combined with Gaifman locality.

\end{proof}

\section{Edge Agreement}

\begin{lemma}

Two cycles with identical signatures agree on all edges
contained in $B_r(v)$ for every vertex $v$.

\end{lemma}

\section{Symmetric Difference}

Let

\[
D=C_1\oplus C_2.
\]

\begin{lemma}

Every component cycle of $D$ has length $\le 2k$.

\end{lemma}

\section{Rank Bound}

Let $M$ denote the cycle incidence matrix.

\[
O=MM^T
\]

is the cycle overlap matrix.

Hence

\[
\mathrm{rank}(O)\le S^L.
\]

\section{Conclusion}

Define

\[
C(k,\Delta,R)
=
\left(
T(k,\Delta,R)\cdot
2^{\Delta^{R+1}}
\right)^{T(k,\Delta,R)(2R+1)}.
\]

Then

\[
\mathrm{COR}(G)\le C(k,\Delta,R).
\]

\end{document}

